\section{机器学习系统的 20 年演化}

\subsection{三大浪潮}

陈天奇将过去 20 年的 ML 系统演化分为三个阶段，每个阶段的系统创新都直接促成了算法的突破：

\textbf{第一阶段：大数据时代（2008--2012）}

标志性事件是 Netflix Challenge 推荐系统比赛，以及今日头条等公司通过推荐系统起家。这个时代的突破来自于数据量从几百个数据点跃升到上亿级别。代表系统包括 Apache Spark、XGBoost，以及李沐和戴文渊在百度做的大规模线性回归（非开源系统）。

\textbf{第二阶段：深度学习 Breakthrough（2012--2017）}

AlexNet 的成功不仅是算法的胜利，更是系统工程的胜利：Alex Krizhevsky 用两块 GPU 训练了一个星期，在当时实验通常只跑几分钟的背景下堪称壮举。Caffe 的出现``非常大程度地推动了 CV 领域 adopt 深度学习''。随后 MXNet、TensorFlow、PyTorch 的出现，让搭建模型从``6 个月写 2 万行 CUDA''变成``几分钟写几行 Python''。

\textbf{第三阶段：大模型时代（2020--至今）}

Transformer 统治了模型架构，注意力转向 scaling law、分布式训练、高效推理。新的系统挑战包括：万卡级别的训练集群、模型并行策略、端侧部署、以及结构化生成（如 XGrammar）。

\begin{table}[H]
\centering
\begin{tabular}{p{2.5cm}p{3cm}p{3cm}p{4cm}}
\toprule
\textbf{阶段} & \textbf{时间} & \textbf{标志性事件} & \textbf{代表系统} \\
\midrule
大数据时代 & 2008--2012 & Netflix Challenge & Spark, XGBoost \\
深度学习 & 2012--2017 & AlexNet, ImageNet & Caffe, MXNet, PyTorch \\
大模型时代 & 2020--至今 & GPT, Transformer & vLLM, Triton, MLC LLM \\
\bottomrule
\end{tabular}
\caption{机器学习系统的三大演化阶段}
\end{table}

\subsection{系统对算法的使能作用}

陈天奇强调，每一次算法突破的背后都有系统层面的革新：

\begin{itemize}
    \item AlexNet 之所以能成功，是因为 Alex Krizhevsky 是``非常强的 hacker''，用两块 GPU 训练了一个星期------在当时实验通常只跑几分钟
    \item Caffe ``非常大程度地推动了 computer vision 领域 adopt 深度学习''
    \item 深度学习框架的出现，让搭建模型从``6 个月写 2 万行 CUDA''变成``几分钟写几行 Python''
\end{itemize}

\begin{importantbox}{算法、数据、算力的三角关系}
``现在大家都知道人工智能依赖于数据、算力和算法。但在我们刚开始做的时候，大部分的关注点都在算法上面。''系统的进步（GPU 加速、分布式训练、自动求导框架）才让算法研究者能快速迭代实验，这种使能关系常常被忽视。
\end{importantbox}

\subsection{学术界的机遇与挑战}

在学术界做 ML 系统研究面临的资源差距是现实的------学校百卡集群 vs 前沿实验室的万卡级别。主持人直接追问：``百卡和万卡还是差两三个数量级。''但陈天奇认为这个差距并没有想象中那么致命：

\begin{itemize}
    \item \textbf{资源约束催生创造力}：``Lack of resource is also an opportunity to make people more creative.'' 这就是为什么他仍然坚持做机器学习编译------因为在学校不可能靠堆 100 个人来解决问题。
    \item \textbf{开源社区合作}：学校不能是一个个体，需要做出能让开源社区共建的工作。像 VRL、Magiline 等成功的项目都有高校因子。
    \item \textbf{垂直突破}：不需要做所有东西，但可以做关键模块做到最好。例如他们做的 XGrammar（结构化数据生成）已成为开源社区做 structured generation 的事实标准之一------agent 输出的 JSON 格式是否正确，这个问题``非常重要，但实际上现在的解决方案都做的不好''。
    \item \textbf{R\&D 层面的差距没那么大}：``很多时候工业界也没有那么富，特别是如果你不是在真正做 pre-training 的团队。'' CMU 最近成立的 Flame Center 有百卡级别的集群，已经足以做很多 kernel 优化和系统研究。
\end{itemize}

\begin{warningbox}{关于``做问题已经被做光了''的焦虑}
陈天奇讲了一个故事：刚开始做机器学习时，戴文渊坐在他旁边听一个关于 SVM 优化的报告，感叹说``如果早两年入行就好了，这些 elegant 的问题似乎都被做光了''。那是 20 年前的事。``It turns out there is always new problem.'' 每个时代都有自己的新问题和新挑战，重要的是能够发现并解决它们。
\end{warningbox}

\subsection{MLSYS 的未来方向}

陈天奇一直是 MLSys Conference 背后的重要推动者。他认为机器学习系统``就像数据库系统一样，是一个问题驱动的方向''。未来几年的重要方向包括：

\begin{enumerate}
    \item \textbf{大模型训练与推理优化}：如何更高效地做 pre-training、fine-tuning 和 serving
    \item \textbf{硬件适配}（Machine Learning Compilation）：硬件的 bring-up 和软件适配仍然是巨大瓶颈
    \item \textbf{AI for Systems}：用 AI Agent 来优化 ML System Engineering 本身
    \item \textbf{Agent 系统基础设施}：结构化生成（如 XGrammar）、tool calling 等
\end{enumerate}

\begin{knowledgebox}{模型趋同会让 MLSYS 衰退吗？}
主持人提出了一个尖锐的问题：大语言模型越来越趋同（Top 5 模型可能占了 99\% 的 inference 量），这会不会让 MLSYS 失去存在价值？陈天奇承认这种可能性不能完全否认，但实际情况更复杂。一方面，硬件本身 ``has to be specialized''------每一代芯片的编程模式都不同。另一方面，``Hardware Lottery'' 意味着当前模型形态可能不是最优的------深度学习之所以繁荣，就是因为百花齐放。MLSYS 不仅是优化当前几种模型，更是 enable 新模型形态的探索。
\end{knowledgebox}

\subsection{本章小结}

机器学习系统的演化始终与算法突破相伴相生。系统不是锦上添花，而是使能技术。学术界在这个领域仍有独特优势：中立性、长期投入的耐心、以及开源社区的天然契合度。陈天奇特别强调，虽然学术界在规模上无法与前沿实验室竞争，但``资源约束催生创造力''------做关键模块做到最好、通过开源社区合作扩大影响力，是学术界的独特路径。

